% Präambel
\documentclass[%
fontsize=12pt,					% Schriftgröße
paper=a4,						% Papierformat
twoside=false, 					% einseitiges (false) oder zweiseitiges (true) Dokument
listof=totoc, 					% Tabellen- und Abbildungsverzeichnis ins Inhaltsverzeichnis
bibliography=totoc,				% Literaturverzeichnis ins Inhaltsverzeichnis aufnehmen
titlepage, 						% Titlepage-Umgebung statt \maketitle
headsepline, 					% horizontale Linie unter Kolumnentitel
%abstracton,					% Überschrift beim Abstract einschalten, Abstract muss dazu in {abstract}-Umgebung stehen
DIV=12,							% Satzspiegeleinstellung, 12 ist Standar bei KOMA
BCOR=6mm,						% Bindekorrektur, die den Seitenspiegel um 3mm nach rechts verschiebt,
cleardoublepage=empty,			% Stil einer leeren eingefügten Seite bei Kapitelwechsel
parskip,							% Absatzabstand bei Absatzwechsel einfügen
ngerman
]{scrbook}			
\documentclass[a4paper,12pt]{article}
\usepackage[utf8]{inputenc}
\usepackage[T1]{fontenc}
\usepackage[ngerman]{babel}
\usepackage{listings}
\usepackage{xcolor}
\usepackage{amsmath}
\usepackage{float}
\usepackage[setspace=false]{scrhack}
\usepackage[utf8]{inputenc} 	% ermöglicht die direkte Eingabe von Umlauten
\usepackage[T1]{fontenc} 		% Ausgabe aller zeichen in einer T1-Codierung (wichtig für die Ausgabe von Umlauten!)
\usepackage{babel} 	% deutsche Trennungsregeln und Übersetzung der festcodierten Überschriften
\renewcaptionname{ngerman}{\contentsname}{Inhaltsverzeichnis}
\renewcaptionname{ngerman}{\bibname}{Literaturverzeichnis}
\setlength{\parindent}{0ex} 	% bei neuem Abschnitt nicht einrücken
%------
% Folgende Einstellungen entsprechen den Vorgaben der Leitlinien
\usepackage[onehalfspacing]{setspace}
% Ende Leitlinien
%------
%------
% Folgende Einstellungen sind bei größeren Arbeiten mit viel Text zu empfehlen.
% Hierbei oben DIV=16 einstellen und Zeile \usepackage[onehalfspacing]{setspace} auskommentieren.
%\linespread{1.2}\selectfont     % Zeilenabstand erhöhen - größere Werte als 1.2 nicht verwenden!!
% Ende Einstellung große Arbeiten mit viel Text.
%------

\newcommand{\lowrule}{%
	\leavevmode \kern.06em\vbox{\hrule width.5em}}

\usepackage{siunitx}			% Vereinfachte Eingabe von Einheiten in Formeln
\sisetup{
	number-unit-product = \;,
	inter-unit-product = \:,
	exponent-product = \cdot,
	output-decimal-marker = {,}
}

\usepackage{graphicx}  			% Einbinden von Grafiken erlauben
\usepackage[format=hang,		% Formatierungen von Unter- / Überschriften
font=normal,
labelfont=bf,
justification=RaggedRight,
singlelinecheck=true,
aboveskip=1mm
]{caption}



\usepackage[backend=biber, %% Hilfsprogramm "biber" beim Compilieren nutzen (statt "biblatex" oder "bibtex")
style=alphabetic, %% Zitierstil (siehe Dokumentation)
natbib=true, %% Bereitstellen von natbib-kompatiblen Zitierkommandos
hyperref=true, %% hyperref-Paket verwenden, um Links zu erstellen
]{biblatex}
\setcounter{biburllcpenalty}{7000}
\setcounter{biburlucpenalty}{8000}
\addbibresource{literature/literaturSimon.bib} %% Einbinden der bib-Datei. Endung .bib unbedingt ergänzen
\addbibresource{literature/literaturPatrik.bib} %% Einbinden mehrerer bib-Dateien mit zusätzlichem \addbibresource - Befehl

\usepackage[normalem]{ulem}
\usepackage{csquotes}

% Folgende Zeilen sind auszukommentieren, falls runde Klammern und ein vgl. bei Zitaten erscheinen sollen.
%\makeatletter
%\renewcommand{\@cite}[2]{(vgl. {#1\if@tempswa , #2\fi})} 
%\renewcommand{\@biblabel}[1]{(#1)}
%\makeatother

\usepackage{pdfpages}

\usepackage{enumitem}			% Erlaubt Änderung der Nummerierung in der Umgebung enumerate

\usepackage{amsmath}			% Ergänzungen für Formeln
\usepackage{textcomp} 			% zum Einsatz von Eurozeichen u. a. Symbolen
\usepackage{eurosym}			% bessere Darstellung Euro-Symbol mit \euro

\usepackage[					% Einstellunge Paket hyperref
hyperfootnotes=false,			% im pfd-Output Fußnoten nicht verlinken
hidelinks						% Entfernen von farbigen Umrandungen der Links
]{hyperref}

\usepackage{makeidx}			% Paket zur Erstellung eines Index
\usepackage[toc]{glossaries}			% Paket zur Erstellung eines Glossars
\newglossaryentry{gls:glossar}
{
	name=Glossar,
	description={Als Glossar wird eine Liste von ausgewählten Begriffen bezeichnet, denen eine jeweilige Erläuterung zugeordnet ist. Ein Glossar wird, sofern gewünscht, im Anhang des Dokuments hinzugefügt.}
}
\usepackage[intoc]{nomencl} 			% zur Erstellung des Abkürzungsberzeichnisses

\usepackage[					% Einstellungen für Fußnoten
bottom,							% Ausrichtung unten
multiple,						% Trennung durch Seperator bei mehreren Fußnoten
hang,
marginal
]{footmisc}

\usepackage{calc}				% Paket zum Berechnen von Längen z.B. 0.8\linewidth

\usepackage{xcolor} 			% einfache Verwendung von Farben in nahezu allen Farbmodellen

\usepackage{listings}			% Darstellung von Quellcode mit den Umgebungen {lstlisting}, \lstinline und \lstinputlisting
\lstset{literate=				% Damit können Umlaute innerhalb Listings geschrieben werden
	{Ö}{{\"O}}1
	{Ä}{{\"A}}1
	{Ü}{{\"U}}1
	{ß}{{\ss}}1
	{ü}{{\"u}}1
	{ä}{{\"a}}1
	{ö}{{\"o}}1
}
\definecolor{mygreen}{rgb}{0,0.6,0}
\definecolor{mygray}{rgb}{0.5,0.5,0.5}
\definecolor{mymauve}{rgb}{0.58,0,0.82}
\lstset{ %
	backgroundcolor=\color{white},   % choose the background color; you must add \usepackage{color} or \usepackage{xcolor}; should come as last argument
	basicstyle=\footnotesize,        % the size of the fonts that are used for the code
	breakatwhitespace=false,         % sets if automatic breaks should only happen at whitespace
	breaklines=true,                 % sets automatic line breaking
	captionpos=t,                    % sets the caption-position to (b) bottom or (t) top
	commentstyle=\color{mygreen},    % comment style
	deletekeywords={...},            % if you want to delete keywords from the given language
	escapeinside={\%*}{*)},          % if you want to add LaTeX within your code
	escapeinside={(*@}{@*)},
	extendedchars=true,              % lets you use non-ASCII characters; for 8-bits encodings only, does not work with UTF-8
	frame=none,	                   	% "single" adds a frame around the code; "none"
	keepspaces=true,                 % keeps spaces in text, useful for keeping indentation of code (possibly needs columns=flexible)
	keywordstyle=\color{blue},       % keyword style
	language=[LaTeX]TeX,             % the language of the code
	morekeywords={*,nomenclature},   % if you want to add more keywords to the set
	numbers=left,                    % where to put the line-numbers; possible values are (none, left, right)
	numbersep=5pt,                   % how far the line-numbers are from the code
	numberstyle=\tiny\color{mygray}, % the style that is used for the line-numbers
	rulecolor=\color{black},         % if not set, the frame-color may be changed on line-breaks within not-black text (e.g. comments (green here))
	showspaces=false,                % show spaces everywhere adding particular underscores; it overrides 'showstringspaces'
	showstringspaces=false,          % underline spaces within strings only
	showtabs=false,                  % show tabs within strings adding particular underscores
	stepnumber=1,                    % the step between two line-numbers. If it's 1, each line will be numbered
	stringstyle=\color{mymauve},     % string literal style
	tabsize=2,	                   % sets default tabsize to 2 spaces
	title=\lstname                   % show the filename of files included with \lstinputlisting; also try caption instead of title
}

% -----------------------------------------------------------------------------------------------------------------
% Zum Aktualisieren des Abkürzungsverzeichnisses (Nomenklatur) bitte auf der Kommandozeile folgenden Befehl aufrufen :
% makeindex <Dateiname>.nlo -s nomencl.ist -o <Dateiname>.nls
% Oder besser: Kann in TexStudio unter Tools-Benutzer als Shortlink angelegt werden
% Konfiguration unter: Optionen-Erzeugen-Benutzerbefehle: makeindex -s nomencl.ist -t %.nlg -o %.nls %.nlo
% -----------------------------------------------------------------------------------------------------------------

% Hier die persönlichen Daten eingeben:

\newcommand{\titel}{Realisierung eines Vier-Gewinnt Roboters}
\newcommand{\untertitel}{ggf. Untertitel mit ergänzenden Hinweisen}
\newcommand{\arbeit}{ Studienarbeit T3\lowrule 3200}
\newcommand{\studiengang}{Elektrotechnik}
\newcommand{\studienrichtung}{Automation}
%\newcommand{\studienschwerpunkt}{}
\newcommand{\autor}{Simon Gschell / Patrik Peters}
\newcommand{\bearbeitungszeitraum}{01.01.2025 - 31.06.2025}
\newcommand{\matrikelnr}{123 456}
\newcommand{\kurs}{TEA22}
%\newcommand{\firma}{Name des Dualen Partners (entfällt ggf. bei Studienarbeit)}
\newcommand{\abgabe}{\today}
\newcommand{\betreuerfirma}{Prof. Dr. ing Thorsten Kever}
%\newcommand{\gutachterdhbw}{Gutachterin / Gutachter der DHBW (nur bei Bachelorarbeit)}

\newcommand{\jahr}{2019}			% für Angabe im Copyright-Vermerk der Titelseite

% Folgende Zeilen definieren Abkürzungen, um Befehle schneller eingeben zu können
\newcommand{\ua}{\mbox{u.\,a.\ }}
\newcommand{\zB}{\mbox{z.\,B.\ }}
\newcommand{\bs}{$\backslash$}
\newcommand*\diff{\mathop{}\!\mathrm{d}}	% Differentialzeichen
\newcommand*\Diff[1]{\mathop{}\!\mathrm{d^#1}} % Differentialzeichen höherer Ableitung
\newcommand*\jj{\mathop{}\!\mathrm{j}}	% Komplexe Zahl j

% Folgende Zeilen weden benötigt, um Tikz und PGF-Plot-Grafiken einzubinden
\usepackage{pgfplots}
\usepackage{pgfplotstable}
\pgfplotsset{compat=newest,width=0.6\linewidth}
\usepgfplotslibrary{smithchart}
\usepackage{tikz}						% Tikz sollte nach Listings Pakete geladen werden.
\usetikzlibrary{arrows}
\usetikzlibrary{positioning}
\usetikzlibrary{calc}
\usetikzlibrary{shapes.geometric, arrows.meta}

\hyphenation{Schrift-ar-ten}

%\BeforeClosingMainAux{% siehe KOMA-Script-Anleitung
%\addcontentsline{toc}{chapter}{\indexname}\stepcounter{page}
%}
%
%\makeindex						% Indexverzeichnis erstellen
\makenomenclature				% Abkürzungsverzeichnis erstellen
\makeglossaries

% -------------------------------------------------------------------------------------------
%                     Beginn des Dokumenteninhalts
% -------------------------------------------------------------------------------------------
\begin{document}
\let\texteuro\euro						% Eingabe \texteuro, € oder \euro erzeugt gleiches Ergebnis
\setcounter{secnumdepth}{3}				% Nummerierungstiefe fürs Inhaltsverzeichnis
\setcounter{tocdepth}{3}
\sffamily								% für die Titelei serifenlose Schrift verwenden

% ------------------------------ Titelei -----------------------------------------------------

\thispagestyle{plain}
\hypersetup{pageanchor=false}
\begin{titlepage}
\enlargethispage{4.0cm}
\sffamily 								% Serifenlose Grundschrift für die Titelseite einstellen

\parbox{0.5\linewidth}{
\begin{flushleft}
% Hier ggf. ein Logo der Firma
\end{flushleft}
}
\parbox{0.5\linewidth}{
\begin{flushright}
	\includegraphics[width=0.4\linewidth]{images/DHBW_d_R_FN_46mm_4c}\\[5ex]
\end{flushright}
}
				

\begin{center}

{\fontsize{20.74pt}{24pt}\selectfont
\textbf{\titel}\\[1.5ex]}
{\fontsize{14pt}{17pt}\selectfont
\textbf{\untertitel}\\[5ex]}
{\fontsize{17pt}{20pt}\selectfont
\textbf{\arbeit}\\[2ex]}
{\fontsize{14pt}{17pt}\selectfont
Studiengang \studiengang\\[2ex]}
{\fontsize{12pt}{14pt}\selectfont
Studienrichtung \studienrichtung\\[1ex]
Duale Hochschule Baden-Württemberg Ravensburg, Campus Friedrichshafen\\[5ex]
von\\[1ex]
\autor\\[15ex]}


\end{center}

\begin{flushleft}
{\fontsize{12pt}{14pt}\selectfont
\begin{tabular}{ll}
Abgabedatum:					& \quad \abgabe \\
Bearbeitungszeitraum:		   	& \quad \bearbeitungszeitraum   \\ 
Matrikelnummer: 			& \quad \matrikelnr \\ 
Kurs: 							& \quad \kurs \\
%Dualer Partner:	 			& \quad \firma \\ % entfällt bei Studienarbeit
Betreuerin / Betreuer:  & \quad \betreuerfirma \\ % Betreuerin / Betreuer der Arbeit
%Gutachterin / Gutachter: & \quad \gutachterdhbw \\ [2ex] % Gutachterin / Gutachter der DHBW (nur bei der Bachelorarbeit erforderlich)
\end{tabular}
}
\end{flushleft}
%%%%% Nachfolgende Zeilen einkommentieren, wenn Copyrightvermerk gewünscht ist
%\begin{flushleft}
%{\fontsize{11pt}{13pt}\selectfont
%Copyrightvermerk:\\
%Dieses Werk einschließlich seiner Teile ist \textbf{urheberrechtlich geschützt}. Jede Verwertung außerhalb der engen Grenzen des Urheberrechtgesetzes ist ohne Zustimmung des Autors unzulässig und strafbar. Das gilt insbesondere für Vervielfältigungen, Übersetzungen, Mikroverfilmungen sowie die Einspeicherung und Verarbeitung in elektronischen Systemen.
%}
%\end{flushleft}
%\begin{flushright}
%{\fontsize{11pt}{13pt}\selectfont \copyright{} \jahr }
%\end{flushright}
\end{titlepage}

\ifthenelse{\boolean{@twoside}}{%
	\cleardoublepage
}{%
	\clearpage
}%

\hypersetup{pageanchor=true}
 				% erzeugt die Titelseite
\pagenumbering{roman}					% kleine, römische Seitenzahlen für Titelei
%% Ggf. folgende Zeile auskommentieren, falls der Sperrvermerk gewünscht ist.
%\chapter*{Sperrvermerk} %*-Variante sorgt dafür, das der Sperrvermerk nicht im Inhaltsverzeichnis auftaucht
%gemäß Ziffer 1.1.14 der Anlage 1 zu §§ 3, 4 und 5  der Studien- und Prüfungsordnung für die Bachelorstudiengänge im Studienbereich Technik der Dualen Hochschule Baden-Württemberg vom 29.09.2017 in der Fassung vom 24.07.2023:
%
%Der Inhalt dieser Arbeit darf weder als Ganzes noch in Auszügen Personen außerhalb des Prüfungsprozesses und des Evaluationsverfahrens zugänglich gemacht werden, sofern keine anders lautende Genehmigung vom Dualen Partner vorliegt.
%
%Musterstadt, den \today \\[4ex]
%
%\rule[-0.2cm]{5cm}{0.5pt} \\
%
%\textsc{\autor} \\[10ex]

\chapter*{Erklärung} %*-Variante sorgt dafür, dass die Erklärung nicht im Inhaltsverzeichnis auftaucht

gemäß Ziffer 1.1.14 der Anlage 1 zu §§ 3, 4 und 5  der Studien- und Prüfungsordnung für die Bachelorstudiengänge im Studienbereich Technik der Dualen Hochschule Baden-Württemberg vom 29.09.2017 in der Fassung vom 24.07.2023.

Wir versichern hiermit, dass wir unsere \arbeit\ mit dem Thema:

\begin{quote}
	\textit{\titel} % -\textit{ \untertitel }
\end{quote}

selbstständig verfasst und keine anderen als die angegebenen Quellen und Hilfsmittel benutzt haben. Wir versicheren zudem, dass die eingereichte elektronische Fassung mit der gedruckten Fassung übereinstimmt.\\[6ex]

Friedrichshafen, den \today \\[1ex]

\rule[-0.2cm]{10cm}{0.5pt} \\

\autor \\[10ex]

\rmfamily

\thispagestyle{empty}

 				% Einbinden der eidestattlichen Erklärung
\chapter*{Kurzfassung} %*-Variante sorgt dafür, das Abstract nicht im Inhaltsverzeichnis auftaucht

In der vorliegenden Studienarbeit wird die Entwicklung eines Roboters für das Spiel „Vier gewinnt“ unter Verwendung des LEGO Spike Prime Systems behandelt. Ziel des Projekts ist es, einen Roboter zu entwerfen, der in der Lage ist, autonom die Rolle eines menschlichen Gegners im Spiel „Vier gewinnt“ zu übernehmen. Der Roboter soll nicht nur die Spielzüge des menschlichen Mitspielers erkennen und darauf reagieren, sondern auch selbstständig Spielsteine in das Spielfeld einwerfen und sicherstellen, dass das Spielfeld für den nächsten Zug bereit ist. Dies erfordert eine präzise Steuerung des Roboters, insbesondere beim Platzieren der Spielsteine, sowie die Fähigkeit, das Spielfeld zu überwachen, um die Position der bereits gesetzten Steine zu erkennen.
Das Projekt umfasst verschiedene technische und organisatorische Aspekte. Dazu gehört die mechanische Konstruktion des Roboters, einschließlich des Mechanismus zum Einwerfen der Spielsteine und die Überwachung des Spielfelds, sowie die Entwicklung der Software, die die Spielzüge und das Verhalten des Roboters steuert. Ein weiterer wichtiger Punkt ist die Sensorik: Der Roboter muss in der Lage sein, die aktuelle Spielsituation durch Farbsensoren zu erfassen, um zu wissen, welche Felder im Spielfeld bereits besetzt sind und wo er seinen nächsten Spielstein platzieren kann.
Zum Abschluss des Projekts wird ein kleines Turnier organisiert, bei dem die von verschiedenen Teams entwickelten Roboterlösungen gegeneinander antreten. Dies bietet die Möglichkeit, den Roboter unter realen Bedingungen zu testen und seine Fähigkeiten im direkten Vergleich mit den Lösungen der anderen Kommilitonen zu messen. Um zusätzliche Motivation zu schaffen, wird die beste Lösung am Ende prämiert, was den Wettbewerbsgedanken fördert und den Anreiz erhöht eine effektive Lösungen zu entwickeln.

\clearpage

\chapter*{Abstract} %*-Variante sorgt dafür, das Abstract nicht im Inhaltsverzeichnis auftaucht
This student research project deals with the development of a robot for the game “Four Wins” using the LEGO Spike Prime system. The aim of the project is to design a robot that is capable of autonomously assuming the role of a human opponent in the game “Four Wins”. The robot should not only recognize the moves of the human opponent and react to them, but also independently place pieces on the playing field and ensure that the playing field is ready for the next move. This requires precise control of the robot, especially when placing the pieces, as well as the ability to monitor the playing field in order to recognize the position of the pieces already placed.
The project involves various technical and organizational aspects. These include the mechanical design of the robot, including the mechanism for inserting the tiles and monitoring the playing field, as well as the development of the software that controls the moves and behavior of the robot. Another important point is the sensor technology: the robot must be able to detect the current game situation using color sensors in order to know which squares in the playing field are already occupied and where it can place its next piece.
At the end of the project, a small tournament is organized in which the robot solutions developed by different teams compete against each other. This offers

Translated with DeepL.com (free version)
\clearpage   			% Einbinden des Abstracts

\tableofcontents						% Erzeugen des Inhalsverzeichnisses

\clearpage

% --------------------------------------------------------------------------------------------
%                    Inhalt der Bachelorarbeit
%---------------------------------------------------------------------------------------------
\pagenumbering{arabic}					% arabische Seitenzahlen für den Hauptteil

\rmfamily

\chapter{Einleitung}
\label{cha:Einleitung}

%Folgende Stichworte können zum Aufbau der Einleitung herangezogen werden.
%
%\begin{itemize}
%\item Hinführung, Begründung, Zweck und Ziel der Aufgabenstellung
%\item Erläuterung der Problemstellung
%\item Konkretisierung der zu lösenden Aufgabe
%\item Gegebenenfalls Formulierung einer Leitfrage oder Forschungsfrage
%\item Ausgangslage, geplante Vorgehensweise, Methoden zur Bearbeitung und Zielsituation
%\item Zum Ende der Einleitung wird eine Kurzübersicht über die Inhalte der Kapitel gegeben: \glqq Die Arbeit ist wie folgt gegliedert: ...\grqq
%\end{itemize}
%
%Die Einleitung wird üblicherweise auf ein bis zwei Seiten als fortlaufender Text geschrieben. Eine weitere Untergliederung in nummerierte Abschnitte ist nicht empfehlenswert, da dies erstens unüblich ist, zweitens die Lesbarkeit nicht begünstigt und drittens die Formulierung der Einleitung erschwert. Weitere Empfehlungen zum Aufbau der Einleitung und des gesamten Dokuments sind z.~B. aus \autocite{DHBW.2021} und \autocite{Lindenlauf.2022} zu entnehmen.
%
%\clearpage
%
%Hinweise: 
%
%\begin{itemize}
%\item Auch in der Einleitung unbedingt zu wichtigen Hintergründen und Fakten Zitate aufführen.\index{Zitat} Zitate bitte in der Form \autocite{Tipler.2019} oder mit Seitenbezug \autocite[66]{Ziegler.2017} oder auch mehrere Zitate  \autocite{Tipler.2019, Ziegler.2017} innerhalb einer eckigen Klammer angeben. Zur besseren Lesbarkeit bitte immer ein Leerzeichen vor dem Zitat einfügen.
%	
%\item Bereits in der Einleitung können Abkürzungen erläutert werden. Grundsätzlich gilt, dass bei der ersten Verwendung einer Abkürzung diese auch erläutert wird. Zum Beispiel können das Antiblockiersystem (ABS) oder die Fahrdynamikregelung (Electronic Stability Control, ESC) als Abkürzungen eingeführt werden. In der Datei \textit{pages/abkuerzungen.tex} sind alle verwendeten Abkürzungen einzufügen. Neben dem verpflichtenden Abkürzungsverzeichnis kann auch ein Glossar hinzugefügt werden. In dieser Vorlage können Glossareinträge in der Datei \textit{pages/glossar.tex} eingefügt werden. Ein \Gls{gls:glossar} ist jedoch nicht verpflichtend.
%\end{itemize}

Strategiespiel sind Spiele, welche nach einer Spieltheorie ablaufen.
$"$Vier-Gewinnt$"$, ist ein Gesellschafts- und Strategiespiel. Das Prinzip des Spiels ist recht einfach. Ein Teilnehmer beginnt, anschließend ist der andere Spieler am Zug.  Gewonne  hat der Teilnhemer, der zuesrt das Ziel erreicht bzw erfüllt hat.
\chapter{Spieltheorie}
\label{cha:Grundlagen}

%
%\section{Einleitung}
%\begin{itemize}
%	\item max. Zeitbegrenzung eines Zuges: -> scannen des Spielfeldes und Rechenzeit müssen begrenzt sein
%\end{itemize}
%\section{Grundlagen}
%\subsection{Geschichte}
%\subsection{Regeln}
%Die Regeln für dieses Strategiespiel sind kinderleicht, was auch ein Argument für die große Beliebtheit bei Jung und Alt ist. Das Spiel besteht aus einem senkrecht stehenden Spielbrett und jeweils aus 21 gelben und roten runde Steine. Diese Steine werden abwechselnd in das hohle Spielbrett mit 42 Aussparungen, sieben Spalten und sechs Reihen, eingeworfen. Der Spieler kann somit eine Spalte auswählen und den Stein fallen lassen. Der eingeworfene Stein besetzt den untersten freien Platz dieser Spalte. Das Spiel wird dann gewonnen, wenn einer der Spieler vier oder mehr Steine in einer waagerechten, senkrechten oder diagonalen Reihe platzieren kann. Der andere Spieler hat somit automatisch verloren. Kommt es zu keiner Bildung einer dieser Kombinationen, endet das Spiel bei Vergabe aller Steine in einem Remis.
%
%\section{spieltheoretische Analyse}
%\subsection{kombinatorisches Spiel}
%Das 4 gewinnt Spiel ist ein kombinatorisches Spiel. Sie bieten einen idealen Rahmen, um Strategien, Entscheidungsfindung und Spieltheorie zu untersuchen. Auch weitere Spiele wie Schach, Dame, Mühle und Tic-Tac-Toe sind ebenfalls kombinatorische Spiele. \\
%Diese folgende fünf Eigenschaften machen ein kombinatorisches Spiel aus:
%\begin{enumerate}
%	\item   \textbf{Zwei-Spieler-Struktur: }Kombinatorische Spiele sind auf zwei Spieler beschränkt, die abwechselnd Züge machen. Dies ermöglicht eine klare Analyse von Strategien und Gegenstrategien.
%	\item 	\textbf{Nullsummencharakter: }Der Gewinn eines Spielers entspricht exakt dem Verlust des anderen. Dies führt zu einer antagonistischen Spielsituation, in der die Interessen der Spieler direkt gegeneinander stehen.
%	\item 	\textbf{Endlichkeit: }Die begrenzte Anzahl von Zugmöglichkeiten und die Unmöglichkeit unendlicher Spielverläufe garantieren, dass jedes Spiel zu einem Abschluss kommt.
%	\item 	\textbf{Perfekte Information: }Beide Spieler haben zu jedem Zeitpunkt vollständige Kenntnis über den Spielstand. Dies eliminiert Unsicherheiten, die in Spielen mit verborgenen Informationen auftreten würden.
%	\item 	\textbf{Determinismus:} Der Ausschluss von Zufallselementen macht das Spiel vollständig vorhersehbar, sofern die Strategien der Spieler bekannt sind.
%\end{enumerate}
%	Durch diese genannten fünf Eigenschaften werden kombinatorische Spiele auch "endliches ZweiPersonen-Nullsummenspiel mit perfekter Information" genannt.
%	
%	\subsection{Zermelos Bestimmtheitssatz}
%	Zermelos Bestimmtheitssatz am Beispiel von „Vier gewinnt“
%	Der Bestimmtheitssatz von Ernst Zermelo zeigt, dass jedes kombinatorische Spiel – wie auch „Vier gewinnt“ – in eine von drei Kategorien eingeordnet werden kann:
%	\begin{enumerate}
%		\item \textbf{Der anziehende Spieler (Spieler 1) hat eine dominante Strategie: }Spieler 1 kann bei optimalem Spiel immer gewinnen.
%		\item \textbf{Der nachziehende Spieler (Spieler 2) hat eine dominante Strategie:} Spieler 2 kann bei optimalem Spiel immer gewinnen.
%		\item \textbf{Keiner der Spieler hat eine dominante Strategie:}Beide Spieler können bei optimalem Spiel ein Unentschieden erzwingen.
%	\end{enumerate}
%	
%	Auf „Vier gewinnt“ angewandt bedeutet dies, dass unabhängig von den individuellen Spielstilen der Spieler das Spiel theoretisch in einer dieser Kategorien liegt. Tatsächlich wurde durch Computersimulationen gezeigt, dass der anziehende Spieler (Spieler 1) bei optimaler Spielweise immer gewinnen kann. Damit gehört „Vier gewinnt“ zur ersten Kategorie.
%	
%	Der Beweis des Bestimmtheitssatzes basiert auf Rückwärtsinduktion. Hierbei wird zuerst die letzte mögliche Zugfolge im Spiel analysiert:
%	
%	Wenn ein Spieler im letzten Zug eine dominante Strategie hat (z. B. durch eine Reihe von vier Steinen gewinnen kann), ist der Gewinn gesichert.
%	Man geht Zug für Zug rückwärts und prüft für jede vorherige Spielsituation, ob ein Spieler seinen Gegner zwingen kann, eine Situation herbeizuführen, in der er selbst gewinnt oder zumindest nicht verliert.
%	Für „Vier gewinnt“ bedeutet das: Spieler 1 kann durch geschicktes Spiel von Beginn an sicherstellen, dass er entweder direkt gewinnt oder Spieler 2 zu Zügen zwingt, die ihm langfristig keine Gewinnoptionen lassen.
%	
%	\textbf{Konsequenzen für „Vier gewinnt“}
%	Aus dem Bestimmtheitssatz folgt, dass „Vier gewinnt“ nicht symmetrisch ist – der erste Spieler hat bei optimaler Spielweise immer einen Vorteil. Dies macht deutlich, dass die Siegchancen in solchen Spielen nicht gleich verteilt sein müssen.
%	
%	Durch die Analyse mit Rückwärtsinduktion wurde festgestellt, dass Spieler 1 mit einer optimalen Strategie den Sieg garantieren kann. Spieler 2 kann nur dann gewinnen, wenn Spieler 1 einen Fehler macht. Somit gehört „Vier gewinnt“ zur ersten Kategorie des Bestimmtheitssatzes.
%
%\section{Strategien und Taktiken}
%
%\section{Algorithmen}
%\subsection{Minimax-Algorithmus}
%\subsection{Alpha-Beta-Algorithums}
%
%\section{Heuristiken}
%
%\section{Lösbarkeit und perfektes Spiel}
%
%\section{Implementierungsaspekte für den LEGO Spike Roboter}
%Die größte Herausforderung bei der Realisierung des Vier-Gewinnt-Roboters besteht in der begrenzten Rechenleistung des verwendeten Mikrocontrollers. Anders als bei leistungsstarken Computern, die in der Lage sind, vollständige Spielalgorithmen zu berechnen und dadurch eine perfekte Strategie zu verfolgen, muss der Roboter mit deutlich eingeschränkten Ressourcen auskommen.
%Ein umfassender Algorithmus, der jede mögliche Zugkombination im Vier-Gewinnt-Spiel analysiert, erfordert erheblichen Speicherplatz und Rechenleistung, da mit jedem neuen Zug die Anzahl der möglichen Spielverläufe exponentiell ansteigt. Auf einem leistungsstarken Computer wäre es theoretisch möglich, eine "perfekte" Strategie zu entwickeln, die den gesamten Spielbaum durchläuft und immer den besten Zug auswählt. Der Mikrocontroller des Roboters hingegen hat nicht die Kapazität, all diese Berechnungen in angemessener Zeit durchzuführen.
%Aus diesem Grund kann der Roboter nur eine begrenzte Anzahl von Zügen im Voraus berechnen. Er muss sich auf Algorithmen stützen, die in der Lage sind, kurzfristige Entscheidungen zu treffen, anstatt langfristige Strategien zu verfolgen. Das bedeutet, dass der Roboter durch Heuristiken (d.h. Daumenregeln oder Annäherungen) gesteuert wird, die ihm helfen, gute, aber nicht immer optimale Züge zu machen. Diese Heuristiken können einfache Prinzipien wie das Verhindern eines unmittelbaren Sieges des Gegners oder das Setzen eigener Spielsteine in strategisch günstige Positionen umfassen.
%\section{Schlussfolgerung}
\chapter{Vorgehen}
\label{cha:vorgehen}

%Je nach Art der Arbeit kann diese Kapitelüberschrift auch \glqq Konzeptentwurf\grqq~lauten.
%
%Beschreibung der Ausgangssituation und des Themenumfelds. Ggf. wird darauf eingegangen, welche Randbedingungen und Einflüsse zu beachten sind.
%
%Anforderungsanalyse und Anforderungsdefinition, nach Möglichkeit strukturiert, um zu einem späteren Zeitpunkt die Anforderungen nachvollziehbar verifizieren zu können.
%
%Herleitung einer Lösung (einer Methodik, eines experimentellen Aufbaus oder von unterschiedlichen Konzepten), Lösungsbewertung und bewusste Wahl des gewählten Vorgehens. An dieser Stelle ist auch auf die Zuverlässigkeit einer Methodik oder auf die Genauigkeit von Untersuchungen einzugehen. Die Überlegungen sollen dazu helfen, mit der angestrebten Lösung die gestellten Anforderungen zu erfüllen, um schließlich die Ziele der Arbeit erreichen zu können.
%
%Bei einer Gegenüberstellung von verschiedenen Lösungsansätzen kann z.~B. eine Nutzwertanalyse helfen. Dabei sind nicht nur z.~B. die Funktion, Leistungsfähigkeit, Umsetzbarkeit und Nutzbarkeit, sondern auch z.~B. wirtschaftliche Aspekte, wie Stück-, Entwicklungskosten oder Ressourcenverbrauch zu berücksichtigen. Sehr bedeutend sind auch Aspekte der Nachhaltigkeit unter Betrachtung des gesamten Lebenszyklus einer erarbeiteten Lösung.
%
%Sowohl bei der Anforderungsdefinition, als auch bei der Lösungsfindung gibt es eine große Anzahl an verschiedenen Methoden. Eine kleine Auswahl ist in der folgenden Aufzählung zu finden.
%
%\begin{itemize}
%\item Anforderungsdefinition mithilfe des Requirements Engineering  \autocite{Pohl.2021}
%\item Systems Engineering Ansatz \autocite{Schluter.2023}
%\item Agile Entwicklungsmethodiken \autocite{Cohn.2010, Martin.2020, Wirdemann.2022}
%\item Klassische Bewertungsverfahren \autocite{Breiing.1997, Zangemeister.2014}
%\end{itemize}
%
%Ziel dieses Kapitels ist, dass auf Basis von umfassend und genau formulierten Anforderungen (ggf. auch Nicht-Zielen) eine Lösungsvielfalt erarbeitet wird, welche anschließend strukturiert bewertet wird, um eine fundierte Begründung für die angestrebte Art der Umsetzung herzuleiten.

\section{Anforderung}
\section{Konzept}
\section{Nutzwertanalyse}
\section{Zeitplan}
\chapter{Umsetzung und Ergebnisse}
\label{cha:umsetzung}

%Je nach Art der Arbeit kann diese Kapitelüberschrift auch \glqq Ergebnisse\grqq~lauten, z.~B. bei rein messtechnischen Aufgaben.
%
%Beschreibung der Umsetzung des zuvor gewählten Vorgehens (theoretische Untersuchung, Erhebungen, Durchführung von Experimenten, Prototypenaufbau, Implementierung eines Prozesses, etc.).
%
%Verifikation anhand der zuvor erarbeiteten Anforderungen und Validierung in Bezug auf das zuvor gestellte Ziel. Diskussion der Ergebnisse. Spätestens hier auch auf die Zuverlässigkeit der gewonnenen Erkenntnisse eingehen (z.~B. anhand der Genauigkeit von Messergebnissen).

\section{Voraussetzungen}
Voraussetzung für dieses Projekt ist ein fundiertes Verständnis der Programmierung sowie Kreativität bei der Konstruktion. Diese Fähigkeiten bringen wir durch unsere abgeschlossene Ausbildung im Bereich Mechatronik und Elektronik mit. Darüber hinaus sind Kenntnisse in der Softwareentwicklung erforderlich, insbesondere im Hinblick auf die Programmierung des LEGO Spike Prime Systems, um den Roboter erfolgreich zu steuern und die Interaktion mit der Spielumgebung zu gewährleisten.
Die Studienarbeit erstreckt sich über zwei Praxisphasen, was einem Zeitraum von insgesamt sechs Monaten entspricht. In dieser Zeit werden die theoretischen Grundlagen, die während des Studiums erlangt haben, in die Praxis umsetzen, um eine vollständige Robotiklösung zu entwickeln, die den Anforderungen des Projekts gerecht wird.

\section{Aufbau}
\section{Software}

\subsection{Herausforderung der Programmierung}
Die größte Herausforderung bei der Realisierung des Vier-Gewinnt-Roboters besteht in der begrenzten Rechenleistung des verwendeten Mikrocontrollers. Anders als bei leistungsstarken Computern, die in der Lage sind, vollständige Spielalgorithmen zu berechnen und dadurch eine perfekte Strategie zu verfolgen, muss der Roboter mit deutlich eingeschränkten Ressourcen auskommen.
Ein umfassender Algorithmus, der jede mögliche Zugkombination im Vier-Gewinnt-Spiel analysiert, erfordert erheblichen Speicherplatz und Rechenleistung, da mit jedem neuen Zug die Anzahl der möglichen Spielverläufe exponentiell ansteigt. Auf einem leistungsstarken Computer wäre es theoretisch möglich, eine "perfekte" Strategie zu entwickeln, die den gesamten Spielbaum durchläuft und immer den besten Zug auswählt. Der Mikrocontroller des Roboters hingegen hat nicht die Kapazität, all diese Berechnungen in angemessener Zeit durchzuführen.
Aus diesem Grund kann der Roboter nur eine begrenzte Anzahl von Zügen im Voraus berechnen. Er muss sich auf Algorithmen stützen, die in der Lage sind, kurzfristige Entscheidungen zu treffen, anstatt langfristige Strategien zu verfolgen. Das bedeutet, dass der Roboter durch Heuristiken (d.h. Daumenregeln oder Annäherungen) gesteuert wird, die ihm helfen, gute, aber nicht immer optimale Züge zu machen. Diese Heuristiken können einfache Prinzipien wie das Verhindern eines unmittelbaren Sieges des Gegners oder das Setzen eigener Spielsteine in strategisch günstige Positionen umfassen.
\chapter{Zusammenfassung}


Die spieltheoretische Analyse von 4Gewinnt hat gezeigt, dass das Spiel trotz seiner scheinbaren Einfachheit eine hohe strategische Tiefe besitzt. Durch die Anwendung von Algorithmen wie Minimax und dessen Optimierung durch den Alpha-Beta-Algorithmus können optimale Entscheidungen getroffen werden, um entweder zu gewinnen oder ein Unentschieden zu erzwingen. Aber auch verschiedene Strategien haben einen Einfluss auf den Spielverlauf.

Der Alpha-Beta-Algorithmus ist hierbei besonders hervorzuheben, da er durch die effiziente Kürzung des Suchbaums eine präzise Analyse in kürzerer Zeit ermöglicht. Diese Eigenschaft ist für die praktische Umsetzung auf ressourcenbeschränkten Systemen wie dem LEGO SPIKE Roboter entscheidend. Der Algorithmus gewährleistet, dass der Roboter in der Lage ist, komplexe Spielsituationen in Echtzeit zu bewerten und optimale Züge zu berechnen.

Mit den gewonnenen Erkenntnissen aus der Spieltheorie und ist nun der Grundstein für den zweiten Teil dieser Studienarbeit gelegt. In diesem Teil wird ein Programm entwickelt, das den Alpha-Beta-Algorithmus implementiert und mit der Hardware des LEGO SPIKE Roboters zusammenarbeitet. 




% ---- Literaturverzeichnis ----------
\interlinepenalty 10000					% Verhindert einen Umbruch mitten in Literatureinträgen
\printbibliography						% Erstellen des Literaturverzeichnisses
\interlinepenalty 0
\clearpage

% -----Ausgabe aller Verzeichnisse ---
\setlength{\parskip}{0.5\baselineskip}

% Alle Abkürzungen, die in der Arbeit verwendet werden. Die Alphabetische Sortierung übernimmt Latex. Nachfolgend sind Beispiele genannt, welche nach Bedarf angepasst, gelöscht oder ergänzt werden können.
% Die Angaben in der eckigen Klammer werden zur Sortierung der Einträge verwendet. Vor allem bei Formelzeichen hat man sonst das Problem, dass diese möglicherweise nicht wie gewünscht sortiert werden.

% Bei den unten stehenden Formelzeichen ist erläutert, wie explizite Sortierschlüssel über den Inhalt der eckigen Klammer angegeben werden.

% Zum Aktualisieren des Abkürzungsverzeichnisses (Nomenklatur) bitte auf der Kommandozeile folgenden Befehl aufrufen :
% makeindex <Dateiname>.nlo -s nomencl.ist -o <Dateiname>.nls
% Oder besser: Kann in TexStudio unter Tools-Benutzer als Shortlink angelegt werden
% Konfiguration unter: Optionen-Erzeugen-Benutzerbefehle: makeindex -s nomencl.ist -t %.nlg -o %.nls %.nlo

% Allgemeine Abkürzungen %%%%%%%%%%%%%%%%%%%%%%%%%%%%
%\nomenclature[Abb]{Abb.}{Abbildung}
%\nomenclature[bzw]{bzw.}{beziehungsweise}
%\nomenclature[DHBW]{DHBW}{Duale Hochschule Baden-Württemberg}
%\nomenclature[ebd]{ebd.}{ebenda}
%\nomenclaturev[etal]{et al.}{at alii}
\nomenclature[etc]{etc.}{et cetera}
%\nomenclature[evtl]{evtl.}{eventuell}
\nomenclature[f]{f.}{folgende Seite}
\nomenclature[ff]{ff.}{fortfolgende Seiten}
%\nomenclature[ggf]{ggf.}{gegebenenfalls}
%\nomenclature[Hrsg]{Hrsg.}{Herausgeber}
%\nomenclature[Tab]{Tab.}{Tabelle}
%\nomenclature[ua]{u. a.}{unter anderem}
%\nomenclature[usw]{usw.}{und so weiter}
\nomenclature[vgl]{vgl.}{vergleiche}
\nomenclature[zB]{z. B.}{zum Beispiel}
%\nomenclaturev[zT]{z. T.}{zum Teil}

% Dateiendungen %%%%%%%%%%%%%%%%%%%%%%%%%%%%%%%%%%%%
\nomenclature[EMF]{EMF}{Enhanced Metafile}
\nomenclature[JPG]{JPG}{Joint Photographic Experts Group}
\nomenclature[KI]{KI}{Künstliche Intelligenz}
\nomenclature[PDF]{PDF}{Portable Document Format}
\nomenclature[PNG]{PNG}{Portable Network Graphics}
%\nomenclature[]{XML}{Extensible Markup Language}

% Abkürzungen von Fachbegriffen %%%%%%%%%%%%%%%%%%%%
\nomenclature[ABS]{ABS}{Antiblockiersystem}
\nomenclature[ESC]{ESC}{Electronic Stability Control, Fahrdynamikregelung}

% Formelzeichen %%%%%%%%%%%%%%%%%%%%%%%%%%%%%%%%%%%%
\nomenclature[a]{$a$}{Beschleunigung}
\nomenclature[F]{$F$}{Kraft}
\nomenclature[m]{$m$}{Masse}
\nomenclature[P]{$P$}{Leistung}
\nomenclature[U]{$U$}{Spannung}
\nomenclature[R]{$R$}{Widerstand}


				% Datei mit allgemeinen Abkürzungen laden
\renewcommand{\nomname}{Verzeichnis verwendeter Formelzeichen und Abkürzungen}
%\addcontentsline{toc}{chapter}{\nomname}
\setlength{\nomlabelwidth}{.20\hsize}
\renewcommand{\nomlabel}[1]{#1 \dotfill}
\setlength{\nomitemsep}{-\parsep}
\printnomenclature						% Erzeugen des Abkürzungsverzeichnises, siehe auch Inhalt der Datei pages/abkuerzungen.tex
\clearpage

\listoffigures 							% Erzeugen des Abbildungsverzeichnisses 
\clearpage

\listoftables 							% Erzeugen des Tabellenverzeichnisses
\clearpage

% -----Anhang ------------------------

%\appendix
%\pagenumbering{Roman}					% große, römische Seitenzahlen für Anhang, falls gewünscht
\addchap{A Nutzung von Künstliche Intelligenz basierten Werkzeugen}
\setcounter{chapter}{1}

Im Rahmen dieser Arbeit wurden Künstliche Intelligenz (KI)\index{Künstliche Intelligenz} basierte Werkzeuge benutzt. Tabelle~\ref{tab:anhang_uebersicht_KI_werkzeuge} gibt eine Übersicht über die verwendeten Werkzeuge und den jeweiligen Einsatzzweck.

\begin{table}[hbt]	
	\centering
	\renewcommand{\arraystretch}{1.5}	% Skaliert die Zeilenhöhe der Tabelle
	\captionabove[Liste der verwendeten Künstliche Intelligenz basierten Werkzeuge]{Liste der verwendeten KI basierten Werkzeuge}
	\label{tab:anhang_uebersicht_KI_werkzeuge}
	\begin{tabular}{>{\raggedright\arraybackslash}p{0.3\linewidth} >{\raggedright\arraybackslash}p{0.65\linewidth}}
		\textbf{Werkzeug} & \textbf{Beschreibung der Nutzung}\\
		\hline 
		\hline
		ChatGPT & 	\vspace{-\topsep}
		\begin{itemize}[noitemsep,topsep=0pt,partopsep=0pt,parsep=0pt] 
			\item Grundlagenrecherche zu Spieltheorie 
			
		\end{itemize} \\
		Perplexity &	\vspace{-\topsep}
		\begin{itemize}[noitemsep,topsep=0pt,partopsep=0pt,parsep=0pt] 
			\item Grundlagenrecherche zu Spieltheorie 
			\item Formulierungshilfe
		\end{itemize} \\ 
		DeepL	&	\vspace{-\topsep}
		\begin{itemize}[noitemsep,topsep=0pt,partopsep=0pt,parsep=0pt] 
			\item Unterstützung beim Übersetzung von Abstract
		\end{itemize} \\ 
		Languagetool &	\vspace{-\topsep}
		\begin{itemize}[noitemsep,topsep=0pt,partopsep=0pt,parsep=0pt] 
			\item Formulierungshilfe 
			\item Rechtschreibkorrektur 
		\end{itemize} \\ 
	
	
		\hline 
	\end{tabular} 
\end{table}
\clearpage

%\		% Zeile auskommentieren bei finalem Dokument!


%\renewcommand{\glossaryname}{Glossar}
%\printglossaries
%\addcontentsline{toc}{chapter}{\glossaryname}

\renewcommand{\indexname}{Sachwortverzeichnis}
%%\printindex								% Erzeugen des Sachwortverzeichnisses

\end{document}